\documentclass{beamer}
\usepackage{amsmath}
\setbeameroption{show notes on second screen=right}
\setbeamertemplate{note page}[plain]
\setbeamerfont{note page}{size=\footnotesize}
\beamertemplatenavigationsymbolsempty

\begin{document}

\begin{frame}
\frametitle{}

What is the point of the y-intercept of a line that has a slope of $-\frac23$ and passes through $(6, 2)$?
\note{\alert<1>{What is the point of the y-intercept of a line that has a slope of -2/3 and passes through (6, 2)?}}

\vskip10pt
\pause
Point-slope form: $y - y_1 = m(x-x_1)$
\note{\alert<2>{Since we have a point and a slope, using the point-slope form y minus y1 equals m times the quantity x minus x1 is fastest. We plug in the coordinates of a point for x1 and y1, and leave x and y themselves alone. We also plug in for the slope where m is.}}
\pause \vskip4pt
$y - 2 = -\frac23(x - 6)$
\note{\alert<3>{In this case, we get y minus 2 equals negative 2 over 3 times the quantity x minus 6. We actually don't need to bother with distributing the -2/3 or any other work of that kind, since the question just wants the y-intercept.}}

\pause\note{\alert<4>{Instead, what we can do is plug in x = 0 to find the y-coordinate of the y-intercept. Why does this work?}}
\vskip10pt
\pause\emph{The y-intercept is the point where the line crosses the y-axis. On the y-axis, the x-coordinate is always 0.}
\note{\alert<5>{The y-intercept is on the y-axis, and all along the y-axis, the x-coordinate (the horizontal coordinate) is 0. To remain on the y-axis, we can't have any movement to the left or right, which is why the horizontal coordinate, the x-coordinate, is 0.}}

\vskip10pt
\pause $y-2 = -\frac23(0-6)$
\note{\alert<6>{When we plug in x = 0, we get y minus 2 equals negative 2 over 3 times the quantity 0 minus 6.}}

\vskip4pt
\pause $y-2 = -\frac23(-6)$
\note{\alert<7>{When we simplify the right side, we get y minus 2 equals negative 2 over 3 times negative 6. }}

\vskip4pt
\pause $y-2 = 4$
\note{\alert<8>{Simplify the right side. Negative two thirds times negative six is four. So y minus 2 equals 4.}}

\vskip4pt
\pause $y = 6$
\note{\alert<9>{Adding 2 to both sides, we get y equals 6. This is the y-coordinate of the y-intercept.}}

\vskip4pt
\pause The point of the $y$-intercept is $(0, 6)$.
\note{\alert<10>{Presented as a point, the point of the y-intercept is (0, 6).}}

\end{frame}



%% Blank frames at the end to prevent weird shifts.
\begin{frame}
\end{frame}
\begin{frame}
\end{frame}
\begin{frame}
\end{frame}
\begin{frame}
\end{frame}
\begin{frame}
\end{frame}
\begin{frame}
\end{frame}

\end{document}
